\documentclass[11pt,a4paper]{article}
\usepackage[utf8]{inputenc}
\usepackage[T1]{fontenc}
\usepackage[spanish,es-tabla]{babel}
\usepackage{amsmath,amssymb,booktabs,graphicx,hyperref,geometry,float}
\usepackage[ruled,vlined,linesnumbered]{algorithm2e}
\geometry{margin=2.5cm}

\title{Problema de Cobertura de Conjuntos (Set Cover):\\
Soluci\'on exacta mediante programaci\'on entera y Algoritmo Gen\'etico}
\author{Juan Jose Vargas Correa \and Juan Diego Chaparro Garcia}
\date{\today}

\begin{document}
\maketitle

\begin{abstract}
Este trabajo aborda el problema de cobertura de conjuntos de costo m\'inimo. Se formula como un
programa lineal entero binario resuelto con el solver CBC y se compara con un Algoritmo Gen\'etico
con operador de reparaci\'on. Ambos m\'etodos se eval\'uan sobre las mismas instancias, analizando
calidad de soluci\'on, tiempo de c\'omputo y comportamiento del algoritmo evolutivo.
\end{abstract}

\section{Introducci\'on y descripci\'on del problema}
Dado un universo $U=\{1,\dots,m\}$ y una colecci\'on de subconjuntos $S_1,\dots,S_n\subseteq U$ con
costos $c_j>0$, el problema consiste en elegir la subcolecci\'on de menor costo cuya uni\'on sea $U$.
Es uno de los 21 problemas NP-completos de Karp \cite{karp1972,garey1979} y aparece en
asignaci\'on de personal, localizaci\'on de servicios y dise\~no de redes \cite{caprara2000}.
% TODO: agregar 1-2 aplicaciones concretas y motivar por qu\'e se compara un m\'etodo exacto con una metaheur\'istica.

\section{Formulaci\'on matem\'atica}
\paragraph{Datos.} $a_{ij}=1$ si el elemento $i$ pertenece a $S_j$ (0 en otro caso); $c_j$ es el costo de $S_j$.
\paragraph{Variables de decisi\'on.}
$x_j\in\{0,1\}$, con $x_j=1$ si se selecciona el subconjunto $S_j$.
\begin{align}
\min\quad & \sum_{j=1}^{n} c_j\,x_j \label{eq:obj}\\
\text{s.a.}\quad & \sum_{j=1}^{n} a_{ij}\,x_j \ge 1, && i=1,\dots,m \label{eq:cover}\\
& x_j\in\{0,1\}, && j=1,\dots,n \label{eq:bin}
\end{align}
La ecuaci\'on~\eqref{eq:obj} minimiza el costo total; las restricciones~\eqref{eq:cover} obligan a que
cada elemento est\'e cubierto por al menos un subconjunto elegido; \eqref{eq:bin} define la naturaleza
binaria de las variables \cite{wolsey2020}.

\section{Metodolog\'ia exacta}
El modelo se implementa en Python con PuLP \cite{mitchell2011pulp} y se resuelve con CBC \cite{cbc},
un solver de ramificaci\'on y corte (\textit{branch-and-cut}). Se reportan estado, valor objetivo,
subconjuntos seleccionados y tiempo. % TODO: describir brevemente relajaci\'on LP y branch-and-bound.

\section{Algoritmo Gen\'etico}
Siguiendo la estructura cl\'asica \cite{holland1975,goldberg1989} y la adaptaci\'on al Set Cover de
Beasley y Chu \cite{beasley1996}:
\begin{itemize}
  \item \textbf{Codificaci\'on:} cromosoma binario $x\in\{0,1\}^n$ (un gen por subconjunto).
  \item \textbf{Fitness:} costo $f(x)=\sum_j c_jx_j$ (se minimiza). Como todo individuo se repara, siempre es factible y no se requieren penalizaciones.
  \item \textbf{Inicializaci\'on:} cada gen vale 1 con probabilidad $p_0=0.2$ y luego se aplica reparaci\'on.
  \item \textbf{Selecci\'on:} torneo de tama\~no $k=3$.
  \item \textbf{Cruce:} uniforme con probabilidad $p_c=0.9$.
  \item \textbf{Mutaci\'on:} inversi\'on de bits con probabilidad $p_m=0.02$ por gen.
  \item \textbf{Reparaci\'on:} (i) \emph{construcci\'on} voraz \cite{chvatal1979}: mientras queden elementos sin cubrir se agrega el subconjunto con menor costo por elemento nuevo cubierto; (ii) \emph{poda}: se eliminan, en orden aleatorio, subconjuntos redundantes.
  \item \textbf{Elitismo:} los 2 mejores individuos pasan intactos.
  \item \textbf{Parada:} 200 generaciones o 50 generaciones consecutivas sin mejora.
\end{itemize}
\begin{algorithm}[H]
\caption{Algoritmo Gen\'etico con reparaci\'on}
Inicializar poblaci\'on $P$ y repararla; evaluar\;
\While{no se cumpla el criterio de parada}{
  $P'\leftarrow$ \'elite de $P$\;
  \While{$|P'|<N$}{
    Seleccionar padres por torneo; cruzar; mutar; reparar el hijo; agregarlo a $P'$\;
  }
  $P\leftarrow P'$; actualizar mejor soluci\'on\;
}
\Return mejor soluci\'on\;
\end{algorithm}
% TODO: justificar la elecci\'on de parametros (idealmente con un peque\~no estudio de sensibilidad).

\section{Implementaci\'on}
C\'odigo en Python: \texttt{setcover.py} (instancias, modelo exacto, GA), \texttt{experiments.py}
(experimentos), \texttt{app.py} (interfaz Streamlit). Dependencias en \texttt{requirements.txt};
las instancias son reproducibles mediante semillas. % TODO: captura de pantalla de la interfaz.

\section{Experimentos computacionales}
Se generaron instancias aleatorias de tama\~nos $(m,n)\in\{(30,20),(60,40),(100,80),(150,120),(200,150)\}$,
densidad 0.10 y costos enteros en $[1,10]$, con 3 instancias por tama\~no. El GA se ejecut\'o 10 veces
por instancia (semillas 0--9) por su naturaleza estoc\'astica. % TODO: indicar hardware y versiones.

\begin{table}[H]
\centering\small
\begin{tabular}{rrrrrrrrr}
\toprule
$m$ & $n$ & $z^*$ & $t_{ex}$ (s) & $z_{GA}^{mejor}$ & $z_{GA}^{prom}$ & $t_{GA}$ (s) & Gap prom (\%) & \'Optimos \\
\midrule
30 & 20 & 57.3 & 0.005 & 57.3 & 57.33 & 0.24 & 0.00 & 3/3 \\
60 & 40 & 53.7 & 0.007 & 53.7 & 53.80 & 0.29 & 0.24 & 3/3 \\
100 & 80 & 56.7 & 0.071 & 56.7 & 57.07 & 0.40 & 0.67 & 3/3 \\
150 & 120 & 40.3 & 0.079 & 40.3 & 40.63 & 0.50 & 0.68 & 3/3 \\
200 & 150 & 37.7 & 0.224 & 37.7 & 37.77 & 0.64 & 0.28 & 3/3 \\
\bottomrule
\end{tabular}
\caption{Promedios por tama\~no. El gap es $100\,(z_{GA}-z^*)/z^*$; ``\'Optimos'' cuenta las instancias donde el mejor de las 10 corridas alcanz\'o $z^*$.}
\label{tab:res}
\end{table}

\section{Comparaci\'on y an\'alisis}
% TODO: completar con tus resultados reales (Tabla~\ref{tab:res}). Puntos que pide la r\'ubrica:
% - Calidad: ¿el GA hall\'o el \'optimo? gap promedio y del mejor.
% - Tiempo: curvas de tiempo vs tama\~no para ambos m\'etodos.
% - Comportamiento del GA: convergencia (figura de la interfaz), variabilidad entre corridas.
% - Ventajas/desventajas: garant\'ia de optimalidad y escalabilidad (exacto) frente a flexibilidad y
%   ausencia de garant\'ias (GA).
% - Limitaciones: instancias aleatorias de tama\~no moderado; CBC resuelve r\'apido estos casos.

\section{Conclusiones}
% TODO

\bibliographystyle{plain}
\bibliography{referencias}
\end{document}
